Estas notas foram elaboradas com o objetivo de acompanhar o curso de Álgebra Avançada, oferecendo uma apresentação organizada e detalhada dos principais conceitos, resultados e técnicas que constituem a base da álgebra moderna. O enfoque adotado busca equilibrar rigor matemático e clareza expositiva, servindo tanto como material de apoio para o estudo quanto como referência futura.

O curso está estruturado em três eixos fundamentais. No primeiro, desenvolve-se a teoria de grupos, incluindo o estudo de subgrupos, subgrupos normais e quocientes, a equação de classes e ações de grupos. São também abordados resultados centrais como os teoremas de Sylow, além de classes importantes de grupos, tais como grupos solúveis, simples e livres.

No segundo eixo, introduz-se a teoria de anéis e módulos, com ênfase em construções fundamentais como soma direta, produto direto e produto tensorial. Estudam-se anéis de polinômios e propriedades estruturais de domínios euclidianos, fatoriais e principais, bem como a teoria de ideais. A teoria de módulos sobre domínios de ideais principais permite aplicações relevantes, em particular à classificação de grupos abelianos finitamente gerados. Finalmente, discute-se a noção de anéis noetherianos e o Teorema da Base de Hilbert.

O terceiro eixo é dedicado à teoria de corpos e extensões, abrangendo extensões algébricas e transcendentes, corpos de decomposição, extensões normais e galoisianas. Estuda-se o grupo de Galois associado a um polinômio e culmina-se com o Teorema Fundamental da Teoria de Galois. Também são consideradas propriedades de corpos finitos e outros tópicos complementares.

Estas notas foram elaboradas com base em referências clássicas da literatura, entre as quais destacam-se as obras de Hungerford, Dummit e Foote, Artin, Lang, entre outros, que constituem pilares no estudo da álgebra moderna.

\section*{Ementa}

\begin{enumerate}
    \item \textbf{Grupos:} subgrupos, subgrupos normais, quocientes; equação de classe, ações de grupos; $p$-grupos finitos e os teoremas de Sylow; grupos solúveis; grupos simples; grupos livres.

    \item \textbf{Anéis e módulos:} a soma, o produto direto e o produto tensorial de módulos; anéis de polinômios e Lema de Gauss, domínios euclideanos, fatoriais e principais; ideais bilaterais e maximais; módulos sobre domínios de ideais principais; elementos e módulos de torção; aplicações a grupos abelianos finitamente gerados; anéis noetherianos e o Teorema da Base de Hilbert.

    \item \textbf{Corpos:} extensões algébricas e transcendentes; corpos de decomposição; extensões normais e galoisianas; o grupo de Galois de um polinômio; Teorema Fundamental da Teoria de Galois; corpos finitos. Outros tópicos.
\end{enumerate}

\section*{Bibliografia}

\begin{enumerate}
    \item Hungerford, T.W. \textit{Algebra}. Springer-Verlag, 2000.
    \item Isaacs, I.M. \textit{Algebra: A Graduate Course}. Brooks/Cole Publishing Company, 1994.
    \item Rotman, J.J. \textit{Advanced Modern Algebra}. AMS, 2010 (second edition).
    \item Dummit, D.; Foote, R. \textit{Abstract Algebra}. Wiley, 2004.
    \item Artin, M. \textit{Algebra}. Pearson, 2nd ed.
    \item Lang, S. \textit{Algebra}. GTM 211.
    \item Lorenz, F. \textit{Algebra}, Vol. 1. Springer-Verlag, 2006.
    \item Herstein, I.N. \textit{Topics in Algebra}.
\end{enumerate}
